\subsection{合成任务覆盖与总线事务核对}
本地 12 个任务改变输入参数流长度，同时保持矩阵维度、输出格式和硬件配置一致。逐任务记录见表\ref{tab:synthetic}。每个任务比较 25344 个数并确认 9522 次输出写命令，既检查正常数值，也约束完整矩阵必须实际经过总线写入。输入字数较少的任务仍有相同输出流量，因而不能按输入流长度同比例推断总任务周期。

\begin{table}[H]\centering\small
\caption{本地合成任务的已记录结果}\label{tab:synthetic}
\begin{tabular}{rrrrrl}\toprule
任务 & 输入字数 & 周期 & 输出写命令 & 核对值数 & 结果\\\midrule
00 & 5883 & 23204 & 9522 & 25344 & 通过 \\
01 & 6761 & 24644 & 9522 & 25344 & 通过 \\
02 & 7603 & 26262 & 9522 & 25344 & 通过 \\
03 & 8617 & 28551 & 9522 & 25344 & 通过 \\
04 & 5835 & 23010 & 9522 & 25344 & 通过 \\
05 & 6767 & 24665 & 9522 & 25344 & 通过 \\
06 & 7592 & 26239 & 9522 & 25344 & 通过 \\
07 & 8612 & 28546 & 9522 & 25344 & 通过 \\
08 & 5883 & 23209 & 9522 & 25344 & 通过 \\
09 & 6768 & 24668 & 9522 & 25344 & 通过 \\
10 & 7597 & 26251 & 9522 & 25344 & 通过 \\
11 & 8598 & 28524 & 9522 & 25344 & 通过 \\
\bottomrule\end{tabular}

\end{table}

\begin{figure}[H]\centering
\begin{tikzpicture}
\begin{axis}[width=.92\linewidth,height=62mm,xlabel={输入流字数（32 位字）},ylabel={完整任务周期},
 xmin=5600,xmax=8900,ymin=22000,ymax=29500,xtick={6000,6500,7000,7500,8000,8500},
 xticklabel style={/pgf/number format/.cd,fixed,1000 sep={}},
 grid=major,grid style={gray!20},scaled ticks=false,
 tick label style={font=\small},label style={font=\small},yticklabel style={/pgf/number format/fixed}]
\addplot[only marks,mark=*,accent,mark size=2pt] table[x=input_words,y=cycles,col sep=comma]{synthetic_tasks.csv};
\end{axis}\end{tikzpicture}
\caption{合成任务的输入字数与周期关系。每点对应一个已完成任务；未作拟合，不将相关关系解释为单一瓶颈的测量。}
\end{figure}

合成任务周期从 23010 到 28551 不等，均值为 25647.75。图中任务分布反映输入编码、装载与计算交叠共同形成的运行结果。这些合成任务不是官方最终性能集，因此不能用其中最短任务宣称吞吐，也不能将它们与厂家模型历史任务逐编号相减作为优化增益。

总线检查至少同时观察命令、响应和物理存储结果：命令与响应守恒防止遗漏事务；目标地址和顺序检查防止写错区域；最终 SRAM 解包逐值比对防止仅有正确计数却数据错误。错误注入另行检查失败任务不发布完成标志。三类检查互相补充，单独一项 PASS 都不足以说明完整任务正确。
